\documentclass[a4paper,12pt,oneside]{article}

\usepackage[T2A]{fontenc}
\usepackage[utf8]{inputenc}
\usepackage[russian]{babel}

\usepackage{lmodern}
\usepackage{microtype}

\usepackage{amsmath,amssymb,amsthm,mathtools}

\usepackage{geometry}
\geometry{top=25mm,left=25mm,right=20mm,bottom=25mm}

\usepackage{enumitem}
\setlist[enumerate]{leftmargin=*, itemsep=0.35\baselineskip, topsep=0.35\baselineskip}
\setlist[itemize]{leftmargin=*, itemsep=0.2\baselineskip, topsep=0.25\baselineskip}

\newlist{alphlist}{enumerate}{1}
\setlist[alphlist]{label=\textbf{\alph*)}, leftmargin=*, align=left}

\newlist{numlist}{enumerate}{1}
\setlist[numlist]{label=\textbf{\arabic*)}, leftmargin=*, align=left}

\newcommand{\Task}[1]{%
  \subsection*{#1}%
  \vspace{-0.2\baselineskip}%
}

\usepackage{booktabs}
\usepackage{float}

\setlength{\parindent}{1.25em}
\setlength{\parskip}{0pt}

\begin{document}

\begin{center}
    {\Large\bfseries Кодирование}\\[2mm]
\end{center}
\vspace{2mm}
\hrule
\vspace{4mm}

\Task{Задание 1. Коды Хаффмана и Шеннона--Фано}
Имеется следующий алфавит и вероятностная модель:

\begin{table}[H]
\centering
\begin{tabular}{cc}
\toprule
Символ & Вероятность \\
\midrule
a & 0.35 \\
b & 0.17 \\
c & 0.17 \\
d & 0.16 \\
e & 0.15 \\
\bottomrule
\end{tabular}
\end{table}

\subsection*{a) Код Шеннона--Фано}
Найти код Шеннона--Фано. Определить среднюю длину кодового слова.

\subsection*{b) Код Хаффмана}
Построить код Хаффмана. Определить среднюю длину кодового слова.

\subsection*{c) Сравнение кодов}
Определить, у какого символа код Шеннона--Фано длиннее, чем код Хаффмана, и объяснить, почему выгоднее выбрать более короткое кодовое слово.

\Task{Задание 2. Код Шеннона--Фано}
Построить код для системы из семи букв $\{A,B,C,D,E,F,G\}$ с вероятностями
$0.1$, $0.2$, $0.05$, $0.3$, $0.05$, $0.15$, $0.15$.

Найти среднее количество разрядов. Декодировать последовательность:
\[
10011101001000111101110101111000.
\]

\Task{Задание 3. Код Хаффмана для сочетаний}
Алфавит из трёх символов: $0.75$, $0.1$, $0.15$. Построить:
\begin{alphlist}
    \item код для отдельных букв;
    \item код для двухбуквенных сочетаний.
\end{alphlist}
Определить средние длины и коэффициенты оптимальности.

\Task{Задание 4. Код Хаффмана для блоков}
Алфавит $\{A,B,C\}$ с вероятностями $p(A)=0.6$, $p(B)=0.3$, $p(C)=0.1$.
Построить коды при кодировании:
\begin{alphlist}
    \item по одной букве;
    \item по две буквы;
    \item по три буквы.
\end{alphlist}
Сравнить эффективность.

\Task{Задание 5. Арифметическое кодирование}
Получено сообщение:
\[
0110100110101.
\]
Известно, что в исходном сообщении было: $3$ A, $1$ C, $1$ D, $1$ Y, $2$ N. Декодировать.

\end{document}
